\section[Réalisation physique discrète sur les voies TPK]{Réalisation physique
discrète : lectures positive et orientée sur les voies TPK, domaine commun et
localisation réversible}
\label{sec:int68-realizacion-fisica-discreta}

La forme spectrale, la réalisation de Hilbert, la section dimensionnelle et le
transport de Cartan ne sont pas des décorations ajoutées à une sortie classique. Ce sont
quatre représentations coordonnées d'une même voie TPK enrichie. Pour
exprimer cette unité causale, il faut d'abord construire leur domaine commun ; il ne
suffit pas d'écrire quatre flèches ayant le même symbole de départ.

\subsection{Domaine commun comme produit fibré}

Soit \(\mathsf P_{\rm TPK}^{\rm enr}\) la catégorie des états TPK préparés
et des voies enregistrées compatibles avec leurs extrémités, leur frontière et leurs restrictions.
Soient
\[
 \mathsf D_{\rm sp}:=\mathsf{SurvPath}_{\rm HMT},\qquad
 \mathsf D_Q:=\operatorname{Dom}(\mathcal Q),\qquad
 \mathsf D_D:=\operatorname{Dom}(\mathfrak D_{\rm HMT}),\qquad
 \mathsf D_C:=\operatorname{Dom}(\rho_C^{\rm disc}),
\]
avec leurs foncteurs d'oubli fidèles
\[
 U_i:\mathsf D_i\longrightarrow\mathsf P_{\rm TPK}^{\rm enr}
 \qquad(i\in\{{\rm sp},Q,D,C\}).
\]
Chaque \(U_i\) n'oublie que la réalisation et conserve l'état et la voie
TPK sous-jacents. Nous définissons
\begin{equation}
 \mathsf P_{\rm CT}^{+}
 :=
 \mathsf D_{\rm sp}
 \times_{\mathsf P_{\rm TPK}^{\rm enr}}\mathsf D_Q
 \times_{\mathsf P_{\rm TPK}^{\rm enr}}\mathsf D_D
 \times_{\mathsf P_{\rm TPK}^{\rm enr}}\mathsf D_C.
 \label{eq:int68-pct-pullback}
\end{equation}
Un objet est donc un état préparé muni des quatre
réalisations sur le même antécédent ; une flèche est une voie enregistrée
dont l'image sous-jacente coïncide dans les quatre facteurs. Si les
présentations ne sont identifiées qu'à isomorphisme près, on utilise le produit
fibré bicatégorique correspondant ; dans la présentation fixée du registre,
on adopte le produit fibré strict de
\eqref{eq:int68-pct-pullback}.

\begin{proposition}[Existence d'un objet dans le domaine commun]
\label{prop:int68-pct-no-vacio}
La catégorie \(\mathsf P_{\rm CT}^{+}\) n'est pas vide. Il n'est pas nécessaire de
particulariser la voie électronique centrale pour le prouver : toute section
coinductive certifiée, équipée de sa flèche identité et de ses données neutres, admet
des relèvements compatibles aux quatre domaines de
\eqref{eq:int68-pct-pullback}.
\end{proposition}

\begin{proof}
Le \cref{thm:generacion-monodromica-conjunta-rev7} fournit, par exemple,
une section survivante préparée \(x_\pi\). Son identité
\(\operatorname{id}_{x_\pi}\) est simultanément : une voie survivante ;
un morphisme de \(\mathsf P_{\rm TPK}\), que \(\mathcal Q\) envoie sur
l'identité bornée de \(\mathcal Q(x_\pi)\) ; la voie vide de signature nulle,
que le \cref{thm:funtor-dimensional-minimo} envoie sur
\((0,0,0,1)\) ; et l'identité du groupoïde, que le
\cref{thm:realizacion-discreta-cartan-rev6} envoie sur l'élément neutre
\((I,0)\) de Poincaré. Par oubli, ces quatre relèvements donnent le même objet
\(x_\pi\) et le même morphisme \(\operatorname{id}_{x_\pi}\). Leur quadruplet
définit, par la propriété du produit fibré, un objet de
\(\mathsf P_{\rm CT}^{+}\).
\end{proof}

\begin{definition}[Adaptateurs au domaine commun]
\label{def:int68-pct-positive}
Les quatre adaptateurs sont les projections canoniques
\[
 \begin{aligned}
 F_{\rm sp}:\mathsf P_{\rm CT}^{+}&\longrightarrow
     \mathsf{SurvPath}_{\rm HMT},\\
 F_Q:\mathsf P_{\rm CT}^{+}&\longrightarrow\operatorname{Dom}(\mathcal Q),\\
 F_D:\mathsf P_{\rm CT}^{+}&\longrightarrow
     \operatorname{Dom}(\mathfrak D_{\rm HMT}),\\
 F_C:\mathsf P_{\rm CT}^{+}&\longrightarrow
     \operatorname{Dom}(\rho_C^{\rm disc}).
 \end{aligned}
\]
En particulier, \(F_{\rm sp}\) conserve la famille compatible des
restrictions finies de la section et de la voie ; il ne reconstruit pas cette famille
à partir de son spectre.
\end{definition}

\begin{lemma}[Fonctorialité des adaptateurs]
\label{lem:int68-adaptadores}
Les quatre adaptateurs préservent les identités et la composition.
\end{lemma}

\begin{proof}
L'identité d'un objet du produit fibré est le quadruplet des identités
sur le même état TPK ; chaque projection renvoie l'identité du facteur.
La composition est définie composante par composante entre quadruplets qui
se projettent sur la même concaténation de voies. Projeter avant ou après
composition donne, par définition, le même résultat.
\end{proof}

\subsection{\(4\) composantes et produit positif}

Les théorèmes de réalisation spectrale, hilbertienne, dimensionnelle et de Cartan
fournissent
\[
 \mathfrak S:\mathsf{SurvPath}_{\rm HMT}
      \longrightarrow\mathbf{SpecRoute},
 \qquad
 \mathcal Q:\operatorname{Dom}(\mathcal Q)
      \longrightarrow\mathsf{Hilb}_{\mathbb R,J},
\]
\[
 \mathfrak D_{\rm HMT}:\operatorname{Dom}(\mathfrak D_{\rm HMT})
      \longrightarrow\mathbf B\mathfrak M_D^+,
 \qquad
 \rho_C^{\rm disc}:\operatorname{Dom}(\rho_C^{\rm disc})
      \longrightarrow\mathbf B G_C,
\]
où
\[
 G_C:=\operatorname{Spin}^{+}(1,3)\ltimes\mathbb R^{1,3}
\]
et
\begin{equation}
 \mathfrak M_D^+
 :=\mathbb N\times\mathbb R_{\geq0}^{2}\times\mathbb R_{>0},
 \qquad
 (n,a,b,\lambda)\odot(n',a',b',\lambda')
 =(n+n',a+a',b+b',\lambda\lambda').
 \label{eq:int68-monoide-dimensional-positivo}
\end{equation}
Ici \(\mathbf BM\) est la catégorie à un objet associée au monoïde
\(M\). Le premier foncteur est celui du
\cref{thm:functor-realizacion-espectral-numero} ; le deuxième est la
réalisation de Hilbert définie dans le chapitre des postulats ; le troisième est
le \cref{thm:funtor-dimensional-minimo} ; et le quatrième est le
\cref{thm:realizacion-discreta-cartan-rev6}.

Nous définissons les restrictions communes
\[
 \mathfrak S_+:=\mathfrak S\circ F_{\rm sp},\qquad
 \mathcal Q_+:=\mathcal Q\circ F_Q,\qquad
 \mathfrak D_+:=\mathfrak D_{\rm HMT}\circ F_D,\qquad
 \rho_{C,+}^{\rm disc}:=\rho_C^{\rm disc}\circ F_C.
\]

\begin{theorem}[Réalisation discrète positive sur le domaine commun]
\label{thm:int68-producto-positivo}
L'application
\begin{equation}
 \mathfrak R_{\rm MD}^{\rm disc,+}
 :=(\mathfrak S_+,\mathcal Q_+,\mathfrak D_+,
       \rho_{C,+}^{\rm disc})
 \label{eq:int68-producto-positivo}
\end{equation}
est un foncteur
\[
 \mathsf P_{\rm CT}^{+}\longrightarrow
 \mathbf{SpecRoute}\times\mathsf{Hilb}_{\mathbb R,J}
 \times\mathbf B\mathfrak M_D^+\times\mathbf BG_C.
\]
On n'affirme pas que toute voie TPK appartienne à \(\mathsf P_{\rm CT}^{+}\).
\end{theorem}

\begin{proof}
Par le \cref{lem:int68-adaptadores}, chaque projection \(F_i\) est un foncteur ;
par les théorèmes cités, il en va de même de
\(\mathfrak S,\mathcal Q,\mathfrak D_{\rm HMT}\) et de
\(\rho_C^{\rm disc}\) dans leurs domaines déclarés. Leurs compositions
préservent donc les identités et la concaténation. La composition dans le codomaine
produit se calcule composante par composante, ce qui prouve
\eqref{eq:int68-producto-positivo}.
\end{proof}

La composante de Hilbert a pour cible générale
\(\mathsf{Hilb}_{\mathbb R,J}\), dont les morphismes sont des opérateurs réels
bornés qui entrelacent les structures complexes. Soit
\(\mathsf P_{\rm CT}^{u}\subseteq\mathsf P_{\rm CT}^{+}\) la sous-catégorie
des flèches \(\gamma:x\to y\) pour lesquelles
\begin{equation}
 \mathcal Q_+(\gamma)^*\mathcal Q_+(\gamma)=I_{\mathcal Q_+(x)},
 \qquad
 \mathcal Q_+(\gamma)\mathcal Q_+(\gamma)^*=I_{\mathcal Q_+(y)}.
 \label{eq:int68-unitariedad-doble}
\end{equation}
Les identités et les produits d'opérateurs unitaires sont encore
unitaires ; cette condition définit donc une sous-catégorie, et ce n'est qu'en elle que
l'on obtient
\begin{equation}
 \mathcal U_{\rm reg}:=
 \mathcal Q_+\big|_{\mathsf P_{\rm CT}^{u}}:
 \mathsf P_{\rm CT}^{u}\longrightarrow
 \mathsf{UHilb}_{\mathbb R,J}.
 \label{eq:int68-restriccion-unitaria}
\end{equation}
L'inversibilité unilatérale ou la conservation de la norme sur un sous-ensemble ne
remplacent pas les deux identités de \eqref{eq:int68-unitariedad-doble}.

La notation à sept entrées se décompose en quatre composantes
typées : \(Q\) est \(\mathcal Q_+\) ;
\(U\) est la restriction \eqref{eq:int68-restriccion-unitaria} ; \(A\) est la
coordonnée d'action de \(\mathfrak D_+\) ; \(C\) est la conjugaison déclarée
dans le sous-domaine où elle entrelace la composante hilbertienne et l'orientation ;
\(\mathcal E\) et \(\mathcal B\) sont des applications de lecture constitutives subordonnées aux
données dimensionnelles ; et \(N_{\rm sp}\) est une propriété quadratique de la
réalisation scindée. Aucune des quatre dernières n'est ajoutée au produit comme
un cinquième foncteur sans domaine ni codomaine.

\subsection{Lectures positive et orientée}

Soit \(e\mapsto e^\vee\) le renversement des arêtes du graphe dirigé
sous-jacent. Dans \(\mathsf P_{\rm CT}^{+}\), \(e^\vee\) est une autre arête et n'est
pas encore déclarée inverse catégorique de \(e\). Un poids
\(w(e)=w(e^\vee)>0\) étant fixé, nous définissons sur une voie
\(\gamma=e_1\cdots e_r\)
\[
 V_+(\gamma):=\sum_{j=1}^r w(e_j),
 \qquad
 V_{\rm or}(\gamma):=
 \sum_{j=1}^r\varepsilon(e_j)w(\underline e_j),
\]
où \(\underline e_j\) est l'arête non orientée et
\(\varepsilon(e^\vee)=-\varepsilon(e)\). La première application de lecture conserve
la variation totale positive ; la seconde, le transport net orienté. Toutes deux sont
additives par concaténation et reçoivent la même voie TPK.

\begin{proposition}[Non-factorisation du lecteur positif par le lecteur orienté]
\label{prop:int68-no-factorizacion}
Supposons que l'identité et le lacet
\(ee^\vee:=e^\vee\circ e\) appartiennent au domaine construit. Il n'existe pas de
fonction \(H\) sur \(\operatorname{im}V_{\rm or}\) telle que
\(V_+=H\circ V_{\rm or}\).
\end{proposition}

\begin{proof}
On a
\[
 V_{\rm or}(\operatorname{id})=0
 =V_{\rm or}(ee^\vee),
 \qquad
 V_+(\operatorname{id})=0,
 \qquad
 V_+(ee^\vee)=2w(e)>0.
\]
Une même entrée \(0\) obligerait \(H(0)\) à prendre deux valeurs distinctes.
\end{proof}

Ce résultat est une non-factorisation concrète dans le domaine construit ; ce
n'est pas un théorème universel sur tout opérateur d'entrelacement enrichi possible.
En particulier, il n'autorise pas à supprimer une application de lecture. L'unité réside dans
l'antécédent TPK commun et non dans une compression horizontale entre ses sorties.

\subsection{Localisation réversible et réalisation orientée}

Soit \(\mathsf P_{{\rm CT},{\rm rev}}^+\) la sous-catégorie engendrée par les
arêtes pour lesquelles les quatre conditions suivantes ont été vérifiées :
\begin{enumerate}
 \item \(\mathfrak S_+(e)\) est un isomorphisme de
       \(\mathbf{SpecRoute}\) et
       \(\mathfrak S_+(e^\vee)=\mathfrak S_+(e)^{-1}\) ;
 \item \(\mathcal Q_+(e)\) satisfait
       \eqref{eq:int68-unitariedad-doble} et
       \(\mathcal Q_+(e^\vee)=\mathcal Q_+(e)^*\) ;
 \item la donnée dimensionnelle orientée provient d'une cochaîne additive avec
       \(d_{\rm or}(e^\vee)=-d_{\rm or}(e)\) et prend ses valeurs dans
       \(\operatorname{Gr}(\mathfrak M_D^+)\) ;
 \item les cocycles de Cartan sont normalisés, additifs et antisymétriques
       sous renversement.
\end{enumerate}
Nous formons le groupoïde
\[
 \Pi_\vee(\mathsf P_{{\rm CT},{\rm rev}}^+)
\]
obtenu à partir de la complétion en groupoïde en imposant en outre
\([e^\vee]=[e]^{-1}\). La composante dimensionnelle orientée prend ses valeurs dans
\begin{equation}
 \operatorname{Gr}(\mathfrak M_D^+)
 \cong\mathbb Z\times\mathbb R^2\times\mathbb R_{>0},
 \label{eq:int68-grothendieck-dimensional}
\end{equation}
avec pour inverse
\((n,a,b,\lambda)^{-1}=(-n,-a,-b,\lambda^{-1})\). Ce lecteur signé n'est pas
la grandeur positive \(V_+\) rebaptisée : ses valeurs sur les arêtes
renversées sont distinctes.

Pour la composante de Cartan, soient \(c_g,c_s\) des cocycles comme au point 4,
soit \(R_4\in\operatorname{Spin}^{+}(1,3)\), soit \(R_4u=u\) et soit
\(\ell_0>0\). Alors
\[
 \rho_C^{\rm disc}(\gamma)
 =\bigl(R_4^{c_g(\gamma)},\ell_0c_s(\gamma)u\bigr).
\]
La loi
\((R,a)(S,b)=(RS,a+Rb)\), l'égalité \(R_4u=u\) et l'additivité des
cocycles donnent
\[
 \rho_C^{\rm disc}(\gamma_2\gamma_1)
 =\rho_C^{\rm disc}(\gamma_2)\rho_C^{\rm disc}(\gamma_1),
 \qquad
 \rho_C^{\rm disc}(\gamma^\vee)
 =\rho_C^{\rm disc}(\gamma)^{-1}.
\]

\begin{theorem}[Réalisation discrète orientée]
\label{thm:int68-realizacion-orientada}
Sous les quatre conditions précédentes, les composantes descendent de manière
unique en un foncteur
\[
 \mathfrak R_{\rm MD}^{\rm disc,or}:
 \Pi_\vee(\mathsf P_{{\rm CT},{\rm rev}}^+)
 \longrightarrow
 \operatorname{Core}(\mathbf{SpecRoute})
 \times\mathsf{UHilb}_{\mathbb R,J}
 \times\mathbf B\operatorname{Gr}(\mathfrak M_D^+)
 \times\mathbf BG_C.
\]
En particulier,
\[
 \mathfrak R_{\rm MD}^{\rm disc,or}(\gamma^{-1})
 =\mathfrak R_{\rm MD}^{\rm disc,or}(\gamma)^{-1}.
\]
\end{theorem}

\begin{proof}
Les conditions 1 et 2 rendent inversibles les images spectrale et
hilbertienne ; la condition 3 prend ses valeurs dans un groupe ; et la composante de Cartan
prend déjà ses valeurs dans \(G_C\). Les quatre respectent la relation génératrice
\([e^\vee]=[e]^{-1}\). Par la propriété universelle de la localisation en
groupoïde, chacune descend de manière unique ; leur produit est le foncteur
indiqué. L'identité d'inversion se vérifie composante par composante.
\end{proof}

L'application de lecture positive \(\mathfrak R_{\rm MD}^{\rm disc,+}\) et l'orientée
\(\mathfrak R_{\rm MD}^{\rm disc,or}\) sont sœurs : aucune n'est définie comme
compression de l'autre. Le passage ultérieur d'une valeur orientée à son module n'est
pas, en général, monoïdal et ne coïncide avec une lecture positive que dans un cône
sans annulation. De même, cette paire ne remplace pas la double projection
fondatrice : l'évaluation archimédienne et la branche
Paley--Witt--Golay--Leech--\(V^\natural\)--Monstre--Moonshine demeurent
des applications de lecture sœurs du même état enrichi. Les quatre composantes
construites dans cette section ne remplacent aucune de ces deux projections.

\subsection{Entrelacement de Cartan et limite lisse}

Soit
\[
 \mathcal A_{\ell_0}
 =\begin{pmatrix}\omega&\ell_0^{-1}e\\0&0\end{pmatrix}
\]
une connexion de Cartan déclarée. L'équation
\begin{equation}
 \operatorname{Hol}_{\mathcal A_{\ell_0}}
   \bigl(\mathfrak R_C(\gamma)\bigr)
 =\rho_C\bigl(\operatorname{Hol}_{\rm TPK}(\gamma)\bigr)
 \label{eq:int68-cuadrado-cartan}
\end{equation}
est exacte dans le modèle discret lorsque les deux holonomies sont définies par le
même produit ordonné d'arêtes.

\begin{criterion}[Promotion à une réalisation lisse]
\label{crit:int68-cartan-suave}
L'équation \eqref{eq:int68-cuadrado-cartan} ne définit un prolongement
lisse que si l'on prouve en outre l'indépendance vis-à-vis de la subdivision, la convergence du
produit ordonné et la récupération locale d'une cotétrade et d'une connexion avec
les domaines et la régularité déclarés.
\end{criterion}

Le terme « physique » désigne ici le fait que les quatre composantes publient
le spectre, la réalisation de l'état, les dimensions et le transport de Cartan sans
perdre leur antécédent TPK\@. Il n'affirme ni unicité, ni équivalence de catégories,
ni couverture de toute voie, ni carte SI automatique, ni théorie des champs complète.

\begin{criterion}[Obstructions aux réalisations positive, unitaire et orientée]
\label{crit:int68-obstrucciones-realizacion}
La réalisation positive n'est pas définie si l'un des quatre
adaptateurs manque, si leurs images ne possèdent pas la même voie TPK sous-jacente ou si une
composante ne préserve pas l'identité et la composition. La restriction unitaire
n'existe pas si l'une quelconque des deux égalités de
\eqref{eq:int68-unitariedad-doble} manque. La réalisation orientée ne descend pas au
groupoïde si une image spectrale n'est pas isomorphe, si une relation localisée
change d'image, si une grandeur positive est traitée comme un inverse, si un
renversement ne change pas le signe des données orientées ou si le cocycle de
Cartan cesse d'être normalisé, additif ou antisymétrique. Il n'existe pas de
prolongement lisse compatible si le résultat dépend de la subdivision ou si
la limite ne retrouve pas la cotétrade et la connexion. Toute utilisation d'une coordonnée
externe pour sélectionner la voie viole le domaine construit.
\end{criterion}

\paragraph{Domaine, codomaine et statut des équations de confluence HMT–MD}
Le produit positif est démontré sur le produit fibré commun et la
réalisation orientée, sur la localisation réversible sous les quatre
conditions explicites. Étendre le domaine à d'autres voies exige de construire leurs
quatre relèvements sur un même antécédent. Le prolongement lisse exige,
en outre, les propriétés de subdivision, de convergence et de récupération locale du
\cref{crit:int68-cartan-suave} ; une vérification finie ne remplace pas ces
propriétés.
